\chapter{Background and Related Work}
\label{ch:background}

\section{Medication Adherence and the ABC Taxonomy}

The ABC taxonomy divides medication adherence into initiation,
implementation, and discontinuation, with persistence as the time between the
first and last dose~\cite{vrijens2012taxonomy}. For software, this means ``did
not take the medication'' can describe several different events. MedAdhere maps
the distinction onto six operational states (Figure~\ref{fig:abc-mapping}).
The states \texttt{taken}, \texttt{delayed}, and \texttt{skipped} describe
implementation around one scheduled dose, and \texttt{stopped} records a
discontinuation statement. The states \texttt{question} and \texttt{unknown}
sit off the phase axis: they describe the kind of reply received, not where the
patient stands in treatment. A second, independent axis records the stated
reason, so ``skipped because I ran out'' needs no combined class for each
state--reason pair. Building on a published conceptual frame instead of an
invented binary label set also follows the EMERGE reporting
guideline~\cite{degeest2018emerge}.

% fig-abc-mapping.tex — ABC taxonomy phases mapped onto the six MedAdhere adherence states.
% Source: chapters/background.tex, Section "Medication Adherence and the ABC Taxonomy"
% (lines 4-23), citing vrijens2012taxonomy and degeest2018emerge. Standalone \input-able.
% Needs the decorations.pathreplacing TikZ library (for the persistence brace) in addition
% to the arrows.meta/positioning/fit/shapes.geometric already loaded by main.tex.
\usetikzlibrary{decorations.pathreplacing,calc}
\begin{figure}[htbp]
  \centering
  \resizebox{\textwidth}{!}{%
  \begin{tikzpicture}[
    font=\small,
    phase/.style={draw, rounded corners, align=center, minimum width=4.4cm,
      minimum height=1cm, fill=medDetFill, draw=medDetLine, font=\bfseries\small},
    state/.style={draw, rounded corners, align=center, text width=4.2cm,
      fill=medStoreFill, draw=medStoreLine, inner sep=6pt},
    crosscut/.style={draw, dashed, rounded corners, align=center,
      text width=4.6cm, fill=medExtFill, draw=medExtLine, inner sep=6pt},
    note/.style={draw, dashed, rounded corners, align=left, fill=medNoteFill, draw=medNoteLine,
      text width=6.4cm, inner sep=6pt},
    txt/.style={align=left, text width=4.6cm, font=\small},
    flow/.style={-{Latex[length=2.2mm]}, thick, medInk},
    link/.style={-{Latex[length=1.8mm]}, thin, medSlate}
  ]
    % --- ABC taxonomy timeline: initiation -> implementation -> discontinuation ---
    \node[phase] (init) at (0,0) {Initiation};
    \node[phase] (impl) at (5.4,0) {Implementation};
    \node[phase] (disc) at (10.8,0) {Discontinuation};

    \draw[flow] (init) -- (impl);
    \draw[flow] (impl) -- (disc);

    \draw[decorate, decoration={brace, amplitude=6pt, raise=5pt}, thick]
      (init.north west) -- (disc.north east)
      node[midway, above=11pt, font=\small\itshape]
      {Persistence -- time from initiation to discontinuation};

    % --- Three implementation states, one discontinuation state, no initiation state ---
    \node[crosscut, below=15mm of init] (noinit) {(no operational state\\currently maps here)};
    \node[state, below=15mm of impl] (implstates)
      {\texttt{taken} -- on time\\\texttt{delayed} -- late, with timing\\\texttt{skipped} -- no dose at all};
    \node[state, below=15mm of disc] (discstates) {\texttt{stopped} -- a discontinuation statement};

    \draw[link] (noinit.north) -- (init.south);
    \draw[link] (implstates.north) -- (impl.south);
    \draw[link] (discstates.north) -- (disc.south);

    % --- Cross-cutting states and the second axis, one row below ---
    \node[crosscut, align=left, text width=9.6cm, anchor=north west]
      (crosscutstates) at ($(noinit.south west |- implstates.south)+(0.1,-0.5)$)
      {\textbf{Not phase-mapped.} These states describe the kind of reply, not
       where the patient stands in treatment.\\[3pt]
       \texttt{question} -- medical guidance, not an adherence claim\\
       \texttt{unknown} -- genuine ambiguity, never a forced guess};
    \node[note, align=left, text width=4.2cm, anchor=north east]
      (reasoncode) at ($(discstates.south east |- implstates.south)+(-0.1,-0.5)$)
      {\textbf{Second axis.} \texttt{reason\_code} records \emph{why}:
       \texttt{ran\_out}, \texttt{side\_effect}, \texttt{forgot}.};
  \end{tikzpicture}}
  \caption{The ESPACOMP ABC adherence taxonomy~\cite{vrijens2012taxonomy,degeest2018emerge}
  mapped onto MedAdhere's six operational adherence states. Three states describe
  implementation around one scheduled dose, one records discontinuation, and two
  (\texttt{question}, \texttt{unknown}) are deliberately left unmapped to any single phase.}
  \label{fig:abc-mapping}
\end{figure}


\section{SMS-Based Adherence Interventions}

Text messaging suits adherence work because delivery is asynchronous, the
interface is familiar, and no smartphone application is needed. A
meta-analysis of chronic-disease SMS interventions found improved adherence,
with heterogeneity in interventions and measurement~\cite{thakkar2016sms}. That
supports SMS as a channel without establishing any one reminder design as
clinically effective. Two-way texting is already deployed at service
scale~\cite{feldacker2024twoway}, but human staff triage the replies. The
closest published analogue used unsupervised methods to recover
medication-related intents from free-text messages on a bidirectional
hypertension-management platform~\cite{davoudi2022intents}. That work finds
intent structure after the fact and without a fixed label contract. MedAdhere
classifies prospectively into a predeclared schema shared by the database, the
dashboard, and the evaluation harness.

\section{LLM Classification of Short Text}

Short SMS replies defeat exact keyword matching because meaning arrives as
slang, abbreviation, negation, or indirect phrasing. Large language models
handle that variation~\cite{singhal2023clinical}, but an unconstrained model
can return any text, which makes it a poor system boundary. MedAdhere therefore
pins a small model to one forced tool schema. This places the classifier in the
family of \emph{constrained} generation, which restricts the output space at
decoding time so that only well-formed outputs are reachable. Published methods
decode against a target grammar~\cite{scholak2021picard,geng2023grammar} or run
a finite-state machine over the vocabulary~\cite{willard2023guided}; a
provider's forced tool call is the hosted form of the same idea. Format
restriction has reduced task performance on some
benchmarks~\cite{tam2024format}, which is one reason this project measures the
contract it ships. Constrained decoding is also not trusted alone. Every
response is parsed through a strict shared schema at the boundary, the
discipline King calls ``parse, don't validate''~\cite{king2019parse}.

The inbound text is also untrusted input that may carry instructions. This is
the indirect prompt-injection setting, in which attacker-controlled content
reaches the model through a data channel instead of the operator's
instruction~\cite{abdelnabi2023indirect}; it has since been formalized and
benchmarked~\cite{liu2024promptinjection}. MedAdhere separates instructions
from data in the way Hines et al.\ call
\emph{spotlighting}~\cite{hines2024spotlighting}. The reply is delimited and
marked as content to classify, not instructions to follow. The input is
truncated, and only the forced tool is permitted. These measures reduce the
opportunity for injection without proving the model immune, so the gold set
includes an adversarial slice and ambiguous replies can resolve to
\texttt{unknown}.

\section{Drug Side-Effect Reference Data}

Medication context comes from four sources. RxNorm normalizes product names to
ingredient identifiers~\cite{rxnorm}. openFDA exposes structured product
labeling, including warnings and adverse reactions~\cite{openfda}. Health
Canada's Drug Product Database provides a market check and the
DIN~\cite{dpd}. OnSIDES extracts adverse-event terms from drug labels and links
them to normalized drug concepts~\cite{tanaka2025onsides}. MedAdhere freezes
these facts at enrollment and uses them as reference data, never as a source of
causal claims.

\section{Serverless Architectures for Event-Driven Systems}

The workload is low-volume and event-shaped, so it maps onto managed services:
API Gateway, EventBridge Scheduler for one-time dose and timeout events, a
patient-keyed FIFO queue, and a Lambda for each narrow handler. These services
deliver \emph{at least once}, not exactly once. A carrier can retry a webhook
and a queue can redeliver a message, so every handler must be
idempotent~\cite{helland2012idempotence}. The design uses two stability
patterns from Nygard by name: \emph{bulkhead} isolation and the \emph{circuit
breaker}~\cite{nygard2018releaseit}. Not every component scales to zero. The
database and NAT gateway stay provisioned while the stack is live, so on-demand
deployment and verified teardown are part of the architecture.
